\section{Convergent and singular weights}
\begin{theorem}[Simultaneous relative-Gram stability]\label{thm:finite}
Suppose $L_j\to L_*$, $H_j\to H_*$ and $\alpha_j\to1$, where eventually
$\alpha_j>0$ and
\begin{equation}\label{eq:relativej}
 L_jL_j^T\succeq\alpha_jL_*L_*^T,
 \qquad H_jH_j^T\succeq\alpha_jH_*H_*^T.
\end{equation}
Then $V(L_j,H_j)\to V(L_*,H_*)$ in $[0,+\infty]$.
If $v_*=V(L_*,H_*)<\infty$, the target minimum is attained with PSD mass.
For any admissible sequence satisfying
\[
 \J_{L_j,H_j}(\mu_j)\le V(L_j,H_j)+\eta_j,
 \qquad\eta_j\ge0,\quad\eta_j\to0,
\]
the actions tend to $v_*$. At finite $v_*$ it has a common weak-star
subsequence converging to a target minimizer, every weak-star subsequence
limit is a PSD target minimizer, and along the whole sequence
\begin{equation}\label{eq:extinction}
 L_*^TA=0\Rightarrow\norm{q_jA}_{\TV,F}\to0,
 \qquad H_*^TA=0\Rightarrow\norm{R_jA}_{\TV,F}\to0.
\end{equation}
The action convergence also holds when $v_*=+\infty$.
\end{theorem}
\begin{proof}
If $v_*<\infty$, Lemma~\ref{lem:compactness} gives a target minimizer
$\mu_*$. For this fixed admissible path, the Gram comparison gives,
eventually,
\[
 V(L_j,H_j)\le\J_{L_j,H_j}(\mu_*)
       \le\alpha_j^{-1}v_*.
\]
Thus $\limsup_jV(L_j,H_j)\le v_*$.

For the lower bound, suppose that
$\ell=\liminf_jV(L_j,H_j)<\infty$. Pass to a subsequence along which
the minima converge to $\ell$ through finite values. For each index choose
an admissible path with action between that minimum and that minimum plus
$1/(j+1)$; such a finite-action path exists by the definition of the infimum.
Their actions converge to $\ell$ and are eventually bounded. Lemma~\ref{lem:compactness} gives
a further weak-star subsequence with admissible limit $\mu$, and
Lemma~\ref{lem:jointlsc} gives
\[
 v_*\le\J_{L_*,H_*}(\mu)\le\ell.
\]
This proves the lower bound when $v_*$ is finite. If $v_*=+\infty$,
the same argument rules out a finite lower limit, so the minima tend to
$+\infty$.

Near-minimizer actions are squeezed between the minima and their
vanishing-gap upper bounds. At finite $v_*$ they are eventually bounded,
so Lemma~\ref{lem:compactness} provides a weak-star subsequence. Every
weak-star cluster is admissible, and Lemma~\ref{lem:jointlsc} bounds its
target action by $v_*$; hence it is a target minimizer with PSD mass.
Finally \eqref{eq:trace} bounds the traces along the whole tail of the
sequence. In \eqref{eq:rates}, the energy and trace factors are therefore
bounded, while $L_j^TA\to0$ or $H_j^TA\to0$ in each forbidden target
direction. This proves \eqref{eq:extinction} along the whole sequence.
\end{proof}

\begin{corollary}[Independent fixed eigenbases]\label{cor:bases}
Let $U\in O(k)$ and $W\in O(n)$ be fixed, and suppose all nonnegative
coefficients in
\[
 L_j=U\diag(\ell_j)U^T,\qquad H_j=W\diag(h_j)W^T
\]
converge to their target coefficients. Then all conclusions of
Theorem~\ref{thm:finite} hold, without a separate comparison assumption.
The bases may differ, and the coefficients need not be monotone.
\end{corollary}
\begin{proof}
For each nonzero target coefficient $a_*$ let its error be
$|a_j^2/a_*^2-1|$, and for zero target coefficients let its error be zero.
Let $\delta_j$ be the sum of these errors over both weights, and set
$\alpha_j=\max(1/2,1-\delta_j)$. Then $\alpha_j\to1$, and eventually
$a_j^2\ge\alpha_ja_*^2$ for every coefficient. The zero target entries
require no restriction. Orthogonal congruence gives \eqref{eq:relativej}.
\end{proof}

\begin{corollary}[Reaction zero and partial reaction]\label{cor:reaction}
Fix $L$. If $H_j\to0$, then
\[
 V(L,H_j)\longrightarrow
 b_L:=\inf\{\J_{L,I_n}(G,q,0):(G,q,0)\in\CE\}.
\]
This includes $b_L=+\infty$. At finite $b_L$, balanced minimizers exist,
near-minimizer actions tend to $b_L$, every weak-star cluster is balanced
and optimal, and $\norm{R_j}_{\TV,F}\to0$ along the entire sequence.
More generally, if $H_j\to H_*$ and eventually
$H_jH_j^T\succeq\alpha_jH_*H_*^T$ with $0<\alpha_j\le1$ and
$\alpha_j\to1$, Theorem~\ref{thm:finite} holds with fixed $L$.
The surviving reaction directions remain in the target problem.
\end{corollary}
\begin{proof}
For $H_*=0$ choose $\alpha_j=1$. Finite action at zero reaction weight
forces $R=0$: with $H=0$, the triples $(0,0,C)$ are admissible for every
matrix $C$, so a nonzero local reaction density makes the support supremum
infinite. For $R=0$, deleting $C$ preserves feasibility and
leaves the support pairing unchanged, while the $C=0$ slice is independent
of $H$. Thus the action is independent of $H$ on these paths, and the target
infimum is exactly $b_L$.
For partial reaction, the fixed transport Gram dominates
$\alpha_jLL^T$. Apply Theorem~\ref{thm:finite} in both cases; use $A=I_n$
for full reaction extinction in the zero case.
\end{proof}

\begin{corollary}[Both zero weights]\label{cor:zero}
If $L_j\to0$ and $H_j\to0$, all conclusions of
Theorem~\ref{thm:finite} hold with $(L_*,H_*)=(0,0)$, irrespective of
rotating eigenbases. The same statement applies to zero transport with a
fixed reaction weight. For a positive-definite fixed reaction weight, this
recovers the zero-transport value and minimizer limit of
\cite[Proposition 4.6]{LiZou}. Theorem~\ref{thm:finite} also covers singular
reaction weights, vanishing-gap near-minimizers, and extinction in each
null direction of the limiting weights.
\end{corollary}
\begin{proof}
The zero target Grams satisfy the comparison with factor one automatically.
For fixed reaction, its Gram comparison is equality.
\end{proof}
